This thesis introduced \textsc{ARIA}, a framework for the automated generation of stochastic models from choreographic specifications. By extending choreographic programming with probabilistic branching and, where applicable, stochastic parameters, the framework supports the connection between high-level coordination modeling and quantitative formal analysis.

The evaluation based on representative case studies, including the Modified ThinkTeam protocol and Peer-to-Peer models, suggests that \textsc{ARIA} is capable of constructing explicit-state Markov models for non-trivial distributed systems. Comparisons with PRISM-based reference models indicate that the generated models are structurally meaningful and, for selected cases, show a good degree of agreement with established tool-supported results. A central technical contribution of this work is the implementation of an on-the-fly exploration engine that maintains a global state representation while resolving concurrent interleavings and local program-counter updates during model construction.

Overall, the thesis demonstrates that choreographic specifications can serve not only as a basis for correct-by-construction distributed implementations, but also as a starting point for the automated construction of stochastic models for quantitative formal analysis.

While the current prototype of \textsc{ARIA} demonstrates the feasibility of choreography-based stochastic model generation, several directions for future work remain:
\begin{itemize}
	\item \textbf{Expressiveness and Implementation Support:} While the choreographic paradigm itself is highly expressive, the current prototype implementation of the \textsc{ARIA} engine is restricted to a specific subset of the grammar. A natural next step would be to eliminate these implementation-specific limitations by expanding the engine's capability to resolve additional constructs, such as collective synchronization and more generalized guarded command patterns. Supporting a broader fragment of the grammar would bridge the gap between the theoretical framework and the tool's practical expressiveness, enabling the modeling of a wider range of complex distributed protocols.
	
	\item \textbf{Implementation Scalability and Optimization:} The performance analysis of our prototype reveals that for large-scale models, practical bottlenecks are increasingly driven by implementation-specific factors, such as JVM memory management and the high I/O overhead of model serialization. Future work should therefore focus on technical optimizations of the \textsc{ARIA} backend, specifically by investigating more memory-efficient internal data structures and compact binary export formats. Furthermore, the integration of algorithmic reduction techniques---such as partial-order or symmetry reduction---directly into the exploration engine could significantly increase the tool's capacity to handle the state-space explosion in complex distributed systems.
	
	\item \textbf{Formal Semantics and Correctness Guarantees:} Another important direction is the formal study of the semantic relationship between source choreographies and the generated stochastic models. Establishing a bisimulation result, or a similar correspondence between the operational semantics of the choreography and the induced transition system, would strengthen the theoretical foundations of the approach.
	
	\item \textbf{Testing and Validation Infrastructure:} As the framework matures, its implementation could be strengthened through a more systematic testing strategy. Unit tests for core components such as parsing, semantic rule application, and state deduplication could be complemented by integration tests covering the full pipeline from choreography specification to exported model. Such a test suite would improve confidence in future language extensions and help detect regressions during further development.
\end{itemize}

Among the directions outlined above, a formal semantic correspondence between the choreographic source language and the generated stochastic models appears particularly important. While scalability and language expressiveness are relevant for the practical development of \textsc{ARIA}, the broader motivation of choreographic programming lies in the close connection between high-level specifications and their correct realizations. In the context of stochastic model generation, this suggests the need to complement empirical evaluation with a more rigorous formal account of the translation. For example, establishing a bisimulation result, or a similar semantic correspondence, would provide stronger justification that the transitions of the generated Markov model faithfully reflect the coordination behavior described by the source choreography. Such a result would not only strengthen the theoretical foundations of \textsc{ARIA}, but also increase confidence in the quantitative analyses carried out on the generated models.